\scanpage{186}
% manuscript position: top
\mspage{102}
% manuscript position: 1–3
envisagée, formée des foncteurs qui commutent aux petites $\underline{\lim}$ (resp à celles de cardinal $\leq\pi$).

% manuscript position: 4–16
On peut donc dire aussi que $\operatorname{Ind}_\pi(C)$ joue le rôle d'une ``enveloppe'' stable par \inserted{petites} $\varinjlim$ de la catégorie \deleted{\unsure{$M$}} $C$, quand les $\varinjlim$ \inserted{de cardinal $\leq\pi$, lesquelles} \deleted{qui} existent déjà dans $C$, sont respectées par le plongement de $C$ dans son enveloppe. C'est un résultat \unsure{nullement|réellement} tautologique qu'il suffit pour cela d'ajouter à $C$ des $\varinjlim$ \uline{grandes devant} $\pi$ d'objets de $\unsure{C|c}$, et que \unsure{plus} précisément, la catégorie enveloppe est réalisée par $\operatorname{Ind}_\pi C$ (\unsure{ou}, au choix et moins pléthoriquement, par $\operatorname{Ind}'_\pi C\subset\widehat C$).

% manuscript position: 17–23
\uline{Démonstration} \inserted{Le 1°) est contenu dans 4.4.13 \unsure{a)|b)} et c). \unsure{Prouvons donc|Prouvons le} 2°) \deleted{\unsure{soit|on}}.} Il est évident que $\alpha$\deleted{\unsure{$(C)$}} est un foncteur pleinement fidèle, dont l'image essentielle engendre strictement $M\equiv\operatorname{Ind}_\pi C$, d'où résulte que le foncteur (4.4.16.2) est pl fid, étant induit par le foncteur pleinement fidèle
% manuscript position: 24
\[
\alpha_{\unsure{\mathcal E|E}}^{*}:\underline{\mathrm{Hom}}(M,\mathcal E)\longrightarrow\underline{\mathrm{Hom}}(\deleted{\unsure{M}}C,\mathcal E)
\]
% manuscript position: foot
% archive sheet number: 185
